\chapter{Introduction}
\label{ch:introduction}

\section{Purpose and scope}

This report describes Yap as it exists on the \code{main} branch of the YapMap repository at commit
\code{9271c80}. It covers the architecture of the mobile client and its backend, the data model, the
flow of data through the system, every user-facing feature, and the rules the server enforces. It is
written for engineers who maintain, review or extend the app, and for readers who need to understand
what the app does with personal data.

Every statement in the report is taken from the source code, the SQL migrations in
\code{supabase/migrations}, the repository documentation and the history of merged pull requests. Where
behaviour depends on configuration that lives outside the repository, such as settings in the Supabase
dashboard, the report says so. Numbers such as radii, time limits and rate limits are quoted from the
code that enforces them; \cref{tab:parameters} collects them in one place.

The report is organised in three parts. \Cref{ch:architecture,ch:data-model,ch:server-interface}
describe the structure of the system: its components, its data and the interface between client and
server. \Cref{ch:flows,ch:features} describe behaviour: the major flows step by step, then each feature
with the screens and functions that implement it. \Cref{ch:security,ch:reliability,ch:performance,ch:configuration,ch:testing,ch:limitations}
cover cross-cutting qualities and operations. A glossary and three reference appendices close the report.

\section{What Yap does}

Yap is a mobile app for iOS and Android that helps people in the same place have a spontaneous
conversation in person. The product line on its welcome screen states the intent: \emph{the app that gets
you off the app}. The app is built around a short, repeated loop:

\begin{enumerate}
  \item \textbf{Go live.} A user taps the large \ui{YAP} button, chooses how long to be visible (30, 60 or
    120 minutes) and may add a note of up to 60 characters, such as where to find them. The phone sends
    its position once, and the user becomes \term{live}.
  \item \textbf{See who is nearby.} While live, the user sees a list of the other live users within about
    100 metres. The list shows a photo, first name and last initial, the note, an introduction and
    interests. It never shows a location or a distance.
  \item \textbf{Say hi.} The user opens one person's profile and sends a \term{say-hi request}. The
    recipient has five minutes to answer. A sender can have only one request waiting at a time.
  \item \textbf{Accept or decline.} The recipient sees the sender's full profile and two buttons of
    equal weight, \ui{Accept} and \ui{Not now}. A declined request and an expired request look the same
    to the sender, who sees \ui{Not this time}.
  \item \textbf{Chat and meet.} An accepted request opens a private one-to-one chat. While they chat, both
    users disappear from everyone else's nearby list. The chat exists so the two can find each other in
    the room.
  \item \textbf{End.} Either user ends the chat and records whether they met. If they are still live, both
    return to the nearby lists.
\end{enumerate}

Alongside the loop, the \ui{Map} tab shows \term{events}: small public plans, such as a coffee or a study
break, that a host places at a chosen spot. Anyone within 5 kilometres of an event can join it, and every
member can read and write its group chat. \Cref{fig:core-loop} shows the loop and the event path.

\begin{figure}[t]
  \centering
  % The core loop and the event path (Figure: core loop).
\begin{tikzpicture}[
    node distance=5.5mm,
    loopbox/.style={box, minimum width=21mm, minimum height=11mm, font=\sffamily\footnotesize},
    sub/.style={font=\sffamily\scriptsize, text=muted},
  ]
  \node[layerlabel, inner sep=0pt] at (-1.05,1.45) {ONE-TO-ONE LOOP (WHILE LIVE)};
  \node[loopbox, accentbox] (live) at (0,0) {Go live\\[-1pt]{\scriptsize 30, 60 or 120 min}};
  \node[loopbox, right=of live] (nearby) {Nearby list\\[-1pt]{\scriptsize live users nearby}};
  \node[loopbox, right=of nearby] (hi) {Say hi\\[-1pt]{\scriptsize one at a time}};
  \node[loopbox, right=of hi] (wait) {Waiting\\[-1pt]{\scriptsize up to 5 min}};
  \node[loopbox, accentbox, right=of wait] (chat) {Chat\\[-1pt]{\scriptsize once accepted}};
  \node[loopbox, right=of chat] (end) {Meet and end\\[-1pt]{\scriptsize met or not}};
  \node[loopbox, below=8mm of wait] (no) {Not this time\\[-1pt]{\scriptsize no reason given}};

  \draw[flow] (live) -- (nearby);
  \draw[flow] (nearby) -- (hi);
  \draw[flow] (hi) -- (wait);
  \draw[flow] (wait) -- (chat);
  \draw[flow] (chat) -- (end);
  \draw[flow] (wait) -- node[edgelabel, right=1pt, fill=none, align=left] {declined\\[-1pt]or expired} (no);
  \draw[flow] (no.west) -| (nearby.south);
  \draw[flow] (end.north) -- ++(0,0.32) -| node[edgelabel, pos=0.25, above=0pt, fill=none] {ended: back on the nearby lists if still live} (nearby.north);

  \node[layerlabel, inner sep=0pt] at (-1.05,-3.05) {MAP EVENTS (NO NEED TO BE LIVE)};
  \node[loopbox, accentbox] (map) at (0,-4.0) {Map tab\\[-1pt]{\scriptsize events in 1 to 50 km}};
  \node[loopbox, right=of map] (pin) {Create an event\\[-1pt]{\scriptsize or open a pin}};
  \node[loopbox, right=of pin] (join) {Join\\[-1pt]{\scriptsize within 5 km}};
  \node[loopbox, accentbox, right=of join] (group) {Group chat\\[-1pt]{\scriptsize all members}};
  \draw[flow] (map) -- (pin);
  \draw[flow] (pin) -- (join);
  \draw[flow] (join) -- (group);
  \node[sub, right=4mm of group, align=left, text width=30mm] {A host joins their own event when creating it.};
\end{tikzpicture}

  \caption{The core loop. The upper path is the one-to-one loop; the lower path is the map event feature,
    which does not require being live.}
  \label{fig:core-loop}
\end{figure}

\section{Design principles}

Five principles recur throughout the code. They explain many of the decisions described later.

\begin{description}
  \item[Locations are never shown.] The server returns people, never positions or distances. A live
    position is stored only while the user is live and is never sent back to any client. The only places
    the app shows are event locations, which hosts choose on purpose (\cref{sec:location-privacy}).
  \item[The server enforces every rule.] The client cannot write to any table except its own profile row
    and its own block list. Every other write goes through a database function that checks who is calling,
    what state they are in, whether the two people are near each other or have blocked each other, and
    whether a rate limit applies. A modified client cannot bypass these checks (\cref{ch:security}).
  \item[One thing at a time.] A sender has at most one request waiting, and a user takes part in at most
    one chat. While a request or chat is on screen, further incoming requests are not shown and expire on
    the server.
  \item[No pressure.] Declining is as easy as accepting, and the sender cannot tell a decline from an
    expiry. The interface uses the phrase \ui{Not this time} for both.
  \item[Safety is two taps away.] Every surface that shows another person carries a menu with
    \ui{Report} and \ui{Block}. A block takes effect on the server at once, in both directions.
\end{description}

\section{Conventions}

The report uses the following conventions.

\begin{itemize}
  \item \code{Monospace} marks identifiers, file paths, SQL and values as they appear in the code.
    Paths are relative to the repository root. \ui{Sans-serif} marks labels the app shows on screen.
  \item A \term{database function} is a Postgres function in the \code{public} schema that the client
    calls through the Supabase API (a remote procedure call, or RPC). A \term{helper} is a function in
    the \code{private} schema, which clients cannot call. The \term{caller} is the signed-in user on whose
    behalf a database function runs; functions read the caller's identity from the session with
    \code{auth.uid()} and never accept it as an argument.
  \item A \term{store action} is an exported function in \code{src/state/store.ts}. An \term{API function}
    is an exported function in \code{src/data/api.ts}.
  \item Distances are in metres unless stated otherwise. The client represents instants as milliseconds
    since the Unix epoch; the database stores \code{timestamptz} values and converts them to
    milliseconds in the JSON it returns.
  \item Terms with a specific meaning in Yap, such as \term{live}, \term{open}, \term{chatting},
    \term{person card} and \term{catch-up}, are defined where they first appear and collected in the
    glossary (\cref{ch:glossary}).
\end{itemize}

\section{Key parameters}

\Cref{tab:parameters} lists the limits and timings that define Yap's behaviour, with the place that
enforces each one. Where the client and the server both apply a limit, the server's check is the one that
cannot be bypassed.

\begingroup\tablefont
\begin{longtable}{@{}L{40mm}L{48mm}L{\dimexpr\linewidth-88mm-4\tabcolsep\relax}@{}}
  \caption{Key parameters and where they are enforced.}
  \label{tab:parameters}\\
  \toprule
  \textsf{\textbf{Parameter}} & \textsf{\textbf{Value}} & \textsf{\textbf{Enforced by}} \\
  \midrule
  \endfirsthead
  \toprule
  \textsf{\textbf{Parameter}} & \textsf{\textbf{Value}} & \textsf{\textbf{Enforced by}} \\
  \midrule
  \endhead
  \bottomrule
  \endfoot
  \multicolumn{3}{@{}l}{\textsf{\textbf{Going live and nearby}}} \\
  Live duration & 30, 60 or 120 minutes & \code{go_green}; \code{DURATIONS} in \code{constants.ts} \\
  Live note length & 0 to 60 characters & check constraint; \code{go_green}, \code{update_note} \\
  Nearby radius & 100\,m plus the caller's reported GPS accuracy, at most 150\,m & \code{nearby}, \code{private.near_each_other} \\
  Nearby list size & at most 50 people & \code{nearby} \\
  Nearby refresh & every 20\,s while live and on screen & \code{NEARBY_REFRESH_MS}, \code{useNearbyPolling} \\
  Location changes & at most 3 in any 10 minutes & \code{go_green} \\
  Movement while live & at most 150\,m plus 10\,km/h since the position was set & \code{go_green} \\
  GPS wobble & moves of 150\,m or less are ignored & \code{go_green} \\
  \midrule
  \multicolumn{3}{@{}l}{\textsf{\textbf{Requests and chats}}} \\
  Answer window & 5 minutes & default of \code{chat_requests.expires_at} \\
  Requests sent & at most 5 in any 10 minutes & \code{send_request} \\
  Requests waiting per sender & 1 & partial unique index; \code{send_request} \\
  Chats per user & 1 & \code{send_request}, \code{respond} \\
  Message length & 1 to 1000 characters & check constraint; both send functions \\
  Messages sent & at most 30 per minute, one-to-one and group combined & \code{send_chat_message}, \code{send_event_message} \\
  Chat history loaded & newest 200 messages & \code{HISTORY_LIMIT} in \code{api.ts} \\
  \midrule
  \multicolumn{3}{@{}l}{\textsf{\textbf{Events}}} \\
  Map radius & 1 to 50\,km, default 10\,km & slider and \code{EVENT_RADIUS_KM}; \code{events_near} caps at 50\,km \\
  Visibility after start & 3 hours & \code{events_near} \\
  Start time & now to 7 days ahead (server allows 5 minutes in the past) & event form; \code{create_event} \\
  Events created & at most 5 per host in any 24 hours & \code{create_event} \\
  Events listed & at most 50 & \code{events_near} \\
  Join distance & within 5\,km & \code{join_event} \\
  Title and description & 1 to 60 and 0 to 200 characters & check constraints; form limits \\
  \midrule
  \multicolumn{3}{@{}l}{\textsf{\textbf{Profiles and photos}}} \\
  Names & 1 to 30 characters each & check constraints; form limits \\
  Introduction & 1 to 300 characters & check constraint; form limit \\
  Interests & 3 to 5 & check constraint; setup and edit forms \\
  Extra photos & 4 to 6 in the app; at most 6 on the server & forms; \code{profiles_photo_paths_max} \\
  Uploaded photo & JPEG, at most 1080\,px wide, quality 0.8 & \code{uploadJpeg} \\
  Avatar copy & JPEG, 320\,px wide & \code{uploadJpeg} \\
  Storage object & at most 5\,MiB; JPEG, PNG or WebP & \code{photos} bucket settings \\
  Signed photo URL & valid 1 hour; reused until 10 minutes remain & \code{SIGNED_URL_SECONDS}, \code{SIGNED_URL_MIN_LEFT_MS} \\
  \midrule
  \multicolumn{3}{@{}l}{\textsf{\textbf{Accounts and safety}}} \\
  Password & at least 8 characters & \code{PASSWORD_MIN_LENGTH}; Supabase Auth setting \\
  Emailed code & 8 digits; resend after 60\,s & email templates; \code{verify.tsx} \\
  Automatic hiding & 3 distinct reporters & \code{report_user} \\
  Report details & at most 500 characters & \code{report_user}; check constraint \\
  Housekeeping & every minute & \code{pg_cron} job \code{yap-expire-stale} \\
\end{longtable}
\endgroup
